\chapter{Security}\label{ch:ch08}

Security in APEX was treated as a scheduled engineering activity rather than an
afterthought. A dedicated review was carried out \emph{before} the system left
development, and every issue it surfaced was reachable in the deployed code path
even though none was ever discovered in production. This chapter documents the
result: the threat model that scopes what APEX defends against, the
authentication stack, the layered authorization model, the input-handling
posture, the sandbox boundary for untrusted code, a mapping of the whole system
onto the OWASP Top~10, and the gaps that remain open on the roadmap.

The guiding principle throughout is \emph{fail closed}: on any error,
ambiguity, or missing precondition, the system denies the action rather than
granting it. A missing header, a malformed token, an unknown role, a resource
the caller does not own, or a critical environment variable left unset all
resolve to a refusal, never to a silent ``allow.'' This is the single most
important adjective in the chapter, and it recurs at every layer.

\section{Threat Model}
\label{sec:threat-model}

APEX is a self-hostable classroom portal whose actors---teachers and students
of a single institution---are known to one another. This shapes the threat
model: the realistic adversary is not a nation-state but an insider probing for
data that is not theirs, hostile content flowing through user-authored fields,
a forged or stolen bearer token, and untrusted program code submitted for
grading. \Cref{fig:threat} names the four adversary classes APEX defends
against and pairs each with the concrete control that mitigates it. Naming the
four classes is also an act of scoping: physical server theft, denial-of-service
at the network edge, and supply-chain compromise of third-party dependencies are
explicitly \emph{out} of scope for a graduation project and are expected to be
handled by the hosting environment.

\begin{figure}[htbp]
\centering
\resizebox{\textwidth}{!}{%
\begin{tikzpicture}[
    node distance=6mm and 34mm,
    every node/.style={font=\sffamily\small},
    adv/.style={draw=apexred, fill=apexred!10, rounded corners=2pt,
        align=center, font=\sffamily\scriptsize, text=apexred!55!black,
        minimum width=48mm, minimum height=15mm},
    ctrl/.style={apexbox, align=center, font=\sffamily\scriptsize,
        minimum width=52mm, minimum height=15mm},
    map/.style={apexarrow},
]

% Column headers
\node[apexlabel, font=\sffamily\bfseries\small] (hA) {Adversary classes};

% Left column: four adversaries (red-tinted)
\node[adv, below=8mm of hA.west, anchor=north west] (a1)
    {\textbf{(1) Unauthenticated}\\network attacker\\{\scriptsize\itshape no token; probes public surface}};
\node[adv, below=of a1] (a2)
    {\textbf{(2) Authenticated student}\\privilege escalation\\{\scriptsize\itshape acts with teacher intent}};
\node[adv, below=of a2] (a3)
    {\textbf{(3) Curious student}\\IDOR / CWE-639\\{\scriptsize\itshape probes others' resources by id}};
\node[adv, below=of a3] (a4)
    {\textbf{(4) Malicious code}\\submitted for execution\\{\scriptsize\itshape hostile program in a submission}};

% Right column: four controls (apexbox), aligned to the adversary rows
\node[ctrl, right=of a1] (c1)
    {\textbf{HS256 JWT + bcrypt(10)}\\authentication\\{\scriptsize\ttfamily middleware.Auth}};
\node[ctrl, right=of a2] (c2)
    {\textbf{Role gate}\\{\scriptsize\ttfamily RequireRole("teacher"/"student")}\\{\scriptsize 403 on wrong role}};
\node[ctrl, right=of a3] (c3)
    {\textbf{Per-resource ownership checks}\\{\scriptsize in handlers (Gate 3)}\\{\scriptsize owner\,/\,enrolled\,/\,attempt owner}};
\node[ctrl, right=of a4] (c4)
    {\textbf{\texttt{isolate} sandbox + Judge0}\\{\scriptsize time / memory limits}\\{\scriptsize 2000\,ms\,/\,262144\,KB per problem}};

% Header over the controls column
\node[apexlabel, font=\sffamily\bfseries\small] at (hA-|c1) {Mitigating controls};

% Mappings: adversary -> control
\draw[map] (a1) -- (c1);
\draw[map] (a2) -- (c2);
\draw[map] (a3) -- (c3);
\draw[map] (a4) -- (c4);

% Residual-gaps note beneath the diagram
\node[apexlabel, draw=apexamber, fill=apexamber!10, rounded corners=2pt,
    align=left, font=\sffamily\scriptsize, text=black,
    below=10mm of a4.south west, anchor=north west,
    text width=110mm] (note)
    {\textbf{Residual gaps (honest):} no rate limiting on the auth endpoints
     (bcrypt cost 10 is the only brake on credential stuffing);
     6-character password minimum. Both are documented, not defended.};

\end{tikzpicture}%
}
\caption{Threat model: four adversary classes and their controls.}
\label{fig:threat}
\end{figure}


The four classes read as follows. The \textbf{unauthenticated network
attacker} holds no token and can only probe the public surface---the login,
signup, password-reset, and Google sign-in endpoints, plus the unauthenticated
\texttt{GET /healthz} liveness probe. The \textbf{authenticated student
attempting privilege escalation} holds a valid token but tries to act with
teacher intent, for example by calling a teacher-only route such as
\texttt{POST /exams}. The \textbf{curious student} holds a valid token and the
correct role but probes \emph{other} users' resources by guessing object
identifiers---the classic Insecure Direct Object Reference (IDOR,
CWE-639~\cite{cwe-639}), which is the thematically central risk for a portal
whose whole value is keeping one student's work private from another. Finally,
\textbf{malicious submitted code} treats the grading pipeline as an execution
primitive: a student's program is, by definition, hostile input, and running it
carelessly is remote code execution by design.

A fifth concern cuts across the model: the \textbf{operator who misconfigures a
deploy}. A self-hosting instructor cannot be relied on to set every environment
variable correctly, so the system must fail closed when a critical variable is
missing or weak. The trust boundary is drawn at the process edge---the browser
is untrusted, the server is trusted, and per-resource ownership re-checks happen
\emph{inside} handlers rather than at the middleware boundary, for reasons
developed in \cref{sec:authz}.

\section{Authentication}
\label{sec:authn}

Authentication answers ``who are you?''. APEX proves identity with a signed JSON
Web Token issued after a successful credential check, and it protects the
long-lived secrets---passwords and reset tokens---with the appropriate
one-way primitives.

\subsection{The fail-closed JWT secret}
\label{sec:jwt-secret}

The signing secret is the linchpin of the entire scheme: anyone who knows it can
forge a token for any user. APEX therefore refuses to boot without a strong one.
The configuration loader validates \texttt{JWT\_SECRET} at startup and calls
\texttt{log.Fatal} (which exits with a non-zero status) if it is absent, is the
historical placeholder, or is shorter than 32 characters.

\begin{lstlisting}[language=Go,caption={Fail-closed validation of the signing secret at startup (\texttt{internal/config/config.go}).},label={lst:jwtsecret}]
jwtSecret := os.Getenv("JWT_SECRET")
if jwtSecret == "" ||
    jwtSecret == "dev-secret-change-in-prod" ||
    len(jwtSecret) < 32 {
    log.Fatal("JWT_SECRET must be set to a strong random value (>=32 chars)")
}
\end{lstlisting}

Three foot-guns are rejected: an unset variable, the exact placeholder string
\texttt{dev-secret-change-in-prod} that once served as a silent default, and any
value under 32 characters---a deliberately conservative bar for HMAC-SHA-256.
Because \texttt{log.Fatal} exits non-zero, no process supervisor such as systemd
or Docker can mask the failure; a deploy with a weak secret stops dead rather
than limping on insecurely. This is the mitigation for the misconfiguring
operator, and it embodies the twelve-factor discipline of drawing configuration
from the environment rather than from source~\cite{twelvefactor}. Secrets never
enter version control: the repository ships \texttt{.env.example}, never
\texttt{.env}.

\subsection{Token issuance and verification}
\label{sec:jwt}

Tokens are HS256 (HMAC with SHA-256), a symmetric scheme in which the same
secret both signs and verifies---appropriate here because a single server issues
and validates its own tokens, with no third party needing a public verification
key~\cite{rfc7519}. On successful signup, login, or Google sign-in the server
issues a token whose claims are \texttt{\{user\_id, email, role, name, exp\}},
with \texttt{exp} set 24 hours from issue.

\begin{lstlisting}[language=Go,caption={Token issuance in \texttt{internal/auth/handler.go}.},label={lst:gentoken}]
func generateToken(user models.User) (string, error) {
    claims := jwt.MapClaims{
        "user_id": user.ID,
        "email":   user.Email,
        "role":    user.Role,
        "name":    user.Name,
        "exp":     time.Now().Add(24 * time.Hour).Unix(),
    }
    token := jwt.NewWithClaims(jwt.SigningMethodHS256, claims)
    return token.SignedString([]byte(cfg.JWTSecret))
}
\end{lstlisting}

The 24-hour \texttt{exp} claim is the only automatic kill-switch for a leaked
token, since there is no server-side revocation list (a limitation
revisited in \cref{sec:gaps}). The payload is signed but not encrypted, so no
secret is placed in it; the claims merely mirror what the verification
middleware reads back.

Verification runs on every protected request in \texttt{middleware.Auth}. It
enforces five fail-closed checks in sequence, the most security-critical of
which is \emph{algorithm pinning}.

\begin{lstlisting}[language=Go,caption={JWT verification middleware with algorithm pinning and safe claim extraction (\texttt{internal/middleware/auth.go}).},label={lst:authmw}]
tokenStr := strings.TrimPrefix(header, "Bearer ")
token, err := jwt.Parse(tokenStr, func(t *jwt.Token) (any, error) {
    if _, ok := t.Method.(*jwt.SigningMethodHMAC); !ok {
        return nil, jwt.ErrSignatureInvalid
    }
    return []byte(secret), nil
})
if err != nil || !token.Valid {
    c.AbortWithStatusJSON(http.StatusUnauthorized, gin.H{"error": "invalid token"})
    return
}

claims, ok := token.Claims.(jwt.MapClaims)
if !ok { /* 401 */ }

userIDFloat, ok := claims["user_id"].(float64)
if !ok { /* 401: invalid user_id in token */ }
\end{lstlisting}

The checks are: (1) the \texttt{Authorization} header must be present and begin
with \texttt{Bearer\ }; (2) the token's declared signing method must be an
\texttt{*jwt.SigningMethodHMAC}---if it is not, the key callback returns
\texttt{jwt.ErrSignatureInvalid} and the token is rejected, which closes the
well-known \texttt{alg=none} forgery and the RS256$\rightarrow$HS256 confusion
attack~\cite{owasp-jwt}; (3) \texttt{jwt.Parse} verifies the signature against
the secret and validates \texttt{exp} (expiry checking is built into the v5
parser); (4) the claims must decode to a \texttt{jwt.MapClaims}; and (5) each
claim is extracted with a comma-ok type assertion, so a malformed claim yields a
clean 401 rather than a panic in the request goroutine. Because JSON numbers
decode to \texttt{float64} in Go, \texttt{user\_id} is read as a
\texttt{float64} and converted to \texttt{uint}. Only after all five pass is the
verified identity stashed on the Gin context with \texttt{c.Set}. A request that
fails this check never reaches a handler; this is Gate~1 of the authorization
model.

\subsection{Password storage with bcrypt}
\label{sec:bcrypt}

Passwords are never stored in plaintext and never returned to clients---the
\texttt{User.PasswordHash} field carries the struct tag \texttt{json:"-"}, so it
is never serialized into any response. Hashing uses bcrypt at
\texttt{bcrypt.DefaultCost}, which is cost~10, meaning $2^{10}=1024$ internal
iterations per hash~\cite{provos1999bcrypt}. bcrypt is deliberately slow and
incorporates a per-hash salt automatically, so two users who chose the same
password receive different stored hashes, defeating precomputed rainbow-table
attacks. The same call hashes on signup and on password reset, and login
verifies with \texttt{bcrypt.CompareHashAndPassword}, which is constant-time
within the bcrypt scheme and therefore does not leak how many characters matched
through timing.

A subtle but deliberate defense is \emph{account-enumeration resistance}: the
login handler returns the identical message---\texttt{"invalid email or
password"}---whether the email is unknown, the stored hash is empty (a
Google-only account with no password set), or the password is simply wrong. An
attacker cannot distinguish ``no such account'' from ``wrong password,'' so the
error text cannot be used to harvest valid email addresses.

\subsection{Google sign-in}
\label{sec:google}

APEX supports federated sign-in through Google Identity Services~\cite{gids},
which layers OpenID Connect~\cite{oidc-core} over OAuth~2.0~\cite{rfc6749}. The
frontend obtains an ID token from Google and POSTs it to \texttt{/auth/google}
as \texttt{\{id\_token, role?\}}. The backend verifies the token server-side
with \texttt{google.golang.org/api/idtoken} against \texttt{GOOGLE\_CLIENT\_ID}
(returning 503 if the client ID is unconfigured, so the feature fails closed
when disabled) and requires both \texttt{email\_verified} to be true and a
non-empty subject before trusting any claim. Account resolution then proceeds in
three steps: match by \texttt{google\_id}; otherwise match by email and
\emph{link} the existing local account by stamping its \texttt{google\_id};
otherwise create a new account. When the role is omitted for a brand-new
account, the server returns \texttt{\{needs\_role: true, \dots\}} and the client
re-submits with the chosen role, so a role is never silently defaulted. A
Google-provisioned account has an empty \texttt{PasswordHash}, which is exactly
why password login on such an account returns the same generic failure rather
than an internal error.

\subsection{Password reset}
\label{sec:reset}

The password-reset flow is a self-contained story in fail-closed design, with
six properties worth naming explicitly. First, the raw token is 32 bytes read
from \texttt{crypto/rand}---a cryptographically secure source---and hex-encoded
to 64 characters, making it unguessable. Second, the database never sees the raw
token: only its SHA-256 hash is stored in the \texttt{token\_hash} column (itself
marked \texttt{json:"-"}). Plain SHA-256 is sufficient here, unlike for
passwords, precisely because the token is already 256 bits of uniform
randomness---there is nothing to dictionary-attack---whereas a human-chosen
password needs bcrypt's deliberate slowness. Third, the token expires one hour
after issue. Fourth, requesting a new link invalidates all prior unused tokens
for that user, so only the newest link works. Fifth, a token is single-use: its
\texttt{used\_at} timestamp is stamped on success and checked on every attempt,
so it cannot be replayed even inside its one-hour window. Sixth, the
\texttt{ForgotPassword} endpoint always returns the same 200 response---``If an
account exists for that email, a reset link has been sent''---regardless of
whether the email exists, extending the enumeration resistance of the login path
to the reset path. If \texttt{SMTP\_HOST} is unset, the reset link is logged to
stdout instead of emailed, which keeps local development functional without a
mail server.

\section{Authorization: The Three-Gate Model}
\label{sec:authz}

Authorization answers ``are you allowed to do \emph{this}?''. APEX enforces it
with three independent gates, each of which fails closed, illustrated in
\cref{fig:authz}. A request must clear all three before a handler mutates any
state.

\begin{figure}[htbp]
\centering
\resizebox{\textwidth}{!}{%
\begin{tikzpicture}[
    node distance=6mm and 12mm,
    every node/.style={font=\sffamily\small},
    gate/.style={apexbox, minimum width=28mm, minimum height=14mm, align=center},
    reject/.style={draw=apexred, fill=apexred!10, rounded corners=2pt,
        align=center, font=\sffamily\scriptsize, minimum height=9mm, minimum width=24mm,
        text=apexred!60!black},
    pass/.style={apexlabel, text=apexgreen!55!black, font=\sffamily\scriptsize\bfseries},
    failedge/.style={-{Latex[length=2.4mm]}, draw=apexred, thick},
    okedge/.style={-{Latex[length=2.4mm]}, draw=apexgreen!70!black, thick},
]

% Ingress
\node[apexext, minimum width=22mm, minimum height=14mm] (req)
    {Incoming\\HTTP request};

% Three gates on a row
\node[gate, right=of req] (g1)
    {\textbf{Gate 1}\\Authentication\\{\scriptsize\itshape valid HS256 JWT?}};
\node[gate, right=of g1]  (g2)
    {\textbf{Gate 2}\\Role check\\{\scriptsize\itshape RequireRole:}\\{\scriptsize\itshape teacher\,/\,student}};
\node[gate, right=of g2]  (g3)
    {\textbf{Gate 3}\\Resource ownership\\{\scriptsize\itshape owns class/exam?}\\{\scriptsize\itshape enrolled\,/\,attempt owner?}};

% Success terminal
\node[apexstore, right=of g3, minimum width=24mm, minimum height=14mm] (h)
    {\textbf{Handler}\\executes};

% Reject terminals below the gates
\node[reject, below=14mm of g1] (r401) {\textbf{401}\\Unauthorized};
\node[reject, below=14mm of g2] (r403b) {\textbf{403}\\Forbidden};
\node[reject, below=14mm of g3] (r403c) {\textbf{403}\\Forbidden};

% Forward pass edges (green)
\draw[apexarrow] (req) -- (g1);
\draw[okedge] (g1) -- node[pass, above]{pass} (g2);
\draw[okedge] (g2) -- node[pass, above]{pass} (g3);
\draw[okedge] (g3) -- node[pass, above]{pass} (h);

% Failure branches (red, downward)
\draw[failedge] (g1) -- node[apexlabel, text=apexred!60!black, right]{invalid / missing} (r401);
\draw[failedge] (g2) -- node[apexlabel, text=apexred!60!black, right]{wrong role} (r403b);
\draw[failedge] (g3) -- node[apexlabel, text=apexred!60!black, right]{not owner} (r403c);

\end{tikzpicture}%
}
\caption{Three-gate authorization model.}
\label{fig:authz}
\end{figure}


\textbf{Gate~1---Authentication.} \texttt{middleware.Auth}
(\cref{lst:authmw}) validates the HS256 token and sets \texttt{user\_id},
\texttt{email}, and \texttt{role} on the context. A missing or invalid token is
a 401, and the request never reaches a handler. This gate establishes
\emph{identity} only.

\textbf{Gate~2---Role.} \texttt{middleware.RequireRole(roles\ldots)} reads the
role from context and admits the request only if it matches one of the allowed
roles; a missing role is a 401 and a wrong role is a 403 ``insufficient
permissions.'' Feature subgroups apply it as an allow-list---\texttt{/exams} is
teacher-only, \texttt{/student/*} is student-only, and
\texttt{PUT /submissions/:id/grade} is teacher-only. APEX has exactly two
roles, \texttt{teacher} and \texttt{student}; there is no admin role. A student
calling a teacher-only route is stopped here at Gate~2.

\textbf{Gate~3---Resource ownership.} Gates 1 and 2 know \emph{that} the caller
is a real teacher, but neither knows \emph{which} exam this is or \emph{who owns
it}. Only the handler---which has just loaded the row---knows. Gate~3 is
therefore an in-handler check that the authenticated user owns or may access the
specific resource, and it is the gate that cannot be removed. It is also the
direct mitigation for the curious-student adversary and the fix for the
IDOR finding (CWE-639~\cite{cwe-639}) described in \cref{sec:findings}.

\begin{lstlisting}[language=Go,caption={Gate~3 ownership check on the run-solution endpoint: teachers must own the problem; students must be enrolled in an assigned class (\texttt{internal/submission/handler.go}).},label={lst:gate3}]
var problem models.Problem
if err := database.DB.Select("id", "exam_id", "teacher_id").
    First(&problem, req.ProblemID).Error; err != nil {
    c.JSON(http.StatusNotFound, gin.H{"error": "problem not found"})
    return
}
userID := c.GetUint("user_id")
role, _ := c.Get("role")
if role == "teacher" {
    if problem.TeacherID != userID {
        c.JSON(http.StatusForbidden, gin.H{"error": "not authorized"})
        return
    }
} else {
    if problem.ExamID == nil { /* 403 */ }
    var assigned int64
    database.DB.Model(&models.ExamClass{}).
        Joins("JOIN class_members ON class_members.class_id = exam_classes.class_id").
        Where("exam_classes.exam_id = ? AND class_members.user_id = ?",
            *problem.ExamID, userID).
        Count(&assigned)
    if assigned == 0 {
        c.JSON(http.StatusForbidden, gin.H{"error": "not authorized"})
        return
    }
}
\end{lstlisting}

The problem is loaded by the request-supplied identifier, then the caller's own
identity---not any request field---drives the decision. A teacher may only run
solutions against problems whose \texttt{teacher\_id} equals their own; a student
may only reach a problem that belongs to an exam assigned to a class in which
they are enrolled, verified by a join across \texttt{exam\_classes} and
\texttt{class\_members}. Zero matches is a 403. The same shape recurs across the
codebase: \texttt{GradeSubmission} checks the parent exam's \texttt{teacher\_id};
\texttt{GetSubmission} lets a teacher read only submissions on exams they own and
a student read only their own submission; \texttt{AssignExam} validates that
every supplied \texttt{class\_id} belongs to the calling teacher; and
\texttt{DeleteAccount} scopes every cascade delete by \texttt{user\_id} or
\texttt{teacher\_id} inside a single transaction that rolls back on any failure.

\subsection{Why ownership checks are not centralized}

A natural instinct is to factor every ownership check into one generic
\texttt{RequireOwnership} middleware. APEX deliberately does \emph{not}, because
the rule genuinely differs per endpoint: an exam handler checks against
\texttt{exam.teacher\_id}, a submission handler against the parent exam's
\texttt{teacher\_id}, and an announcement handler against
\texttt{class.teacher\_id} or class membership. Hiding these behind a generic
helper makes the security-relevant rule less visible at the call site, which is
the opposite of what such a check wants. The chosen trade-off is readable
repetition over an abstraction that would be harder to audit---repetition one
can read beats abstraction one cannot verify.

\section{Input Handling}
\label{sec:input}

\subsection{SQL injection}

Every database access goes through GORM's parameterized helpers~\cite{gorm},
which bind values as query parameters rather than interpolating them into SQL
text---structurally closing SQL injection (CWE-89~\cite{cwe-89}). The clause
\texttt{Where("email = ?", req.Email)} sends \texttt{req.Email} as a bound
parameter, so a payload such as \texttt{'~OR~1=1--} is treated as a literal
string, never as SQL. The only raw \texttt{db.Exec} calls in the codebase live
in \texttt{internal/database/migrations.go}, where the SQL is static and
developer-authored, with no user input reaching it.

\subsection{Cross-site scripting}

The API returns JSON only---responses are produced by \texttt{c.JSON}, and there
is no path that assembles an HTML document server-side, so reflected XSS through
the API has no vector. On the client, React escapes text by default, and APEX
never applies \texttt{dangerouslySetInnerHTML} to user-authored content such as
announcement bodies, problem descriptions, MCQ option text, or exam titles. The
one place raw markup could once execute was the Markdown renderer: the
\texttt{rehype-raw} plugin re-parsed inline HTML inside a Markdown source,
turning a problem description into a stored-XSS vector (CWE-79~\cite{cwe-79}).
That plugin was removed; Markdown still renders through \texttt{react-markdown}
with \texttt{remark-gfm} for tables and task lists, but raw HTML is now stripped
at parse time. The trade-off---teachers can no longer hand-write raw HTML in a
description---is acceptable, as Markdown is expressive enough for the intended
content.

\subsection{CSRF and the localStorage trade-off}

Cross-Site Request Forgery (CWE-352~\cite{cwe-352}) exploits the browser's
automatic attachment of ambient credentials---chiefly cookies---to
cross-origin requests. APEX carries no such ambient credential: the JWT is stored
in \texttt{localStorage} under the key \texttt{kernel-token} and is attached
\emph{explicitly} by application code as an \texttt{Authorization:\ Bearer}
header on each request. A cross-site form or image cannot read
\texttt{localStorage} or set an \texttt{Authorization} header, so there is no
classic CSRF surface, and CORS (with \texttt{AllowCredentials:\ false}) further
signals that no credentialed cross-origin access is expected.

This choice is a deliberate trade-off rather than a free win. Storing the token in
\texttt{localStorage} trades CSRF resistance for XSS exposure: a successful XSS
would let injected script read \texttt{localStorage} and exfiltrate the token,
whereas an \texttt{httpOnly} cookie would be unreadable to script but would
reintroduce CSRF concerns and demand anti-CSRF tokens plus same-site controls.
The design was chosen for a simpler self-hosted deployment---a bearer header is
trivial to reason about behind a same-origin proxy---and the residual XSS risk is
narrowed by the \texttt{rehype-raw} removal above, by the absence of
\texttt{dangerouslySetInnerHTML}, and by the 24-hour token expiry that bounds the
blast radius of any stolen token. Hardening the design with \texttt{httpOnly},
same-site cookies and a refresh-token model is named explicitly in
\cref{sec:gaps}.

\section{Sandbox Isolation for Untrusted Code}
\label{sec:sandbox}

Student submissions are arbitrary programs in one of six languages---Python,
JavaScript, TypeScript, Java, C, and C++---and are, by definition, hostile
input. Executing them in the APEX process would be remote code execution by
design: a submission could delete files, read other students' data, open network
connections, spin an infinite loop, or exhaust memory. APEX avoids this entirely
by never executing student code itself. The \texttt{internal/judge0} package is
purely an HTTP client: it marshals \texttt{\{language\_id, source\_code, stdin\}}
to JSON, POSTs it to Judge0~\cite{judge0}, and reads a JSON verdict back, with a
30-second client timeout. There is no \texttt{os/exec}, no \texttt{eval}, and no
shell anywhere in the package; an unsupported language short-circuits to a
sentinel status and never reaches Judge0 at all.

Judge0 in turn runs each submission inside \texttt{isolate}~\cite{isolate}, a
sandbox built on Linux namespaces and cgroups---the same technology used to grade
the International Olympiad in Informatics. Each run receives a throwaway
filesystem that cannot read the host, no network access, a hard CPU-time limit,
and a memory cap. APEX passes the per-problem \texttt{time\_limit\_ms} (default
2000\,ms) and \texttt{memory\_limit\_kb} (default 262144\,KB, i.e.\ 256\,MB) as
those bounds, so an infinite loop becomes a ``time limit exceeded'' verdict
rather than a hung server, and a memory bomb is killed rather than allowed to
crash the host. The most security-sensitive component in the system is thus one
that APEX delegates entirely rather than builds---the safest code to run is code
one does not run. The operational cost of this delegation is a production
dependency on a self-hosted Judge0, since the default public CE instance at
\texttt{https://ce.judge0.com} is rate-limited and unsuitable for real use.

\section{Findings and OWASP Mapping}
\label{sec:findings}

The scheduled pre-ship review found four issues, all reachable in the deployed
code path and all fixed in a single commit before the system left development.
\Cref{tab:findings} summarizes them. The unifying observation is that impact is
not proportional to line count: each fix is a few lines, and each closes a real,
reachable exploit path.

\begin{table}[htbp]
\centering
\caption{Security findings from the pre-ship review, with severity, weakness
class, and fix.}
\label{tab:findings}
\begin{tabularx}{\linewidth}{@{}l l l X@{}}
\toprule
\textbf{ID} & \textbf{Severity} & \textbf{Class} & \textbf{Finding and fix} \\
\midrule
F1 & HIGH & IDOR (CWE-639) & \texttt{POST /submissions/run} trusted the
request-supplied \texttt{problem\_id} and leaked any problem's sample tests.
Fixed by the Gate~3 ownership check (\cref{lst:gate3}). \\
\addlinespace
F2 & HIGH & Stored XSS (CWE-79) & Markdown rendered with \texttt{rehype-raw}
executed embedded raw HTML. Fixed by removing the plugin so raw HTML is stripped
at parse time. \\
\addlinespace
F3 & MED & Weak secret & \texttt{config.Load} defaulted \texttt{JWT\_SECRET} to a
publicly-known placeholder. Fixed by the fail-closed validator
(\cref{lst:jwtsecret}). \\
\addlinespace
F4 & MED & HTML injection & Notification-email bodies interpolated unescaped
user fields. Fixed by passing them through \texttt{html.EscapeString}. \\
\bottomrule
\end{tabularx}
\end{table}

Mapped onto the OWASP Top~10~\cite{owasp-top10} and cross-checked against the
Application Security Verification Standard~\cite{owasp-asvs}, the posture is
summarized in \cref{tab:owasp}. The two categories that dominate a
multi-tenant classroom portal---broken access control and injection---are the
ones the design invests most heavily in.

\begin{table}[htbp]
\centering
\caption{APEX posture against the OWASP Top~10 categories.}
\label{tab:owasp}
\begin{tabularx}{\linewidth}{@{}l X@{}}
\toprule
\textbf{Category} & \textbf{APEX posture} \\
\midrule
A01 Broken Access Control & Three-gate model; Gate~3 per-resource ownership
checks in every handler (fixed F1/IDOR). \\
\addlinespace
A02 Cryptographic Failures & bcrypt cost~10 for passwords; SHA-256-hashed,
single-use reset tokens; HS256 tokens signed with a validated $\geq$32-char
secret. \\
\addlinespace
A03 Injection & GORM parameterized queries (SQLi); JSON-only API and raw-HTML
stripping (XSS, fixed F2). \\
\addlinespace
A05 Security Misconfiguration & Fail-closed \texttt{JWT\_SECRET} validation
(fixed F3); secrets kept out of source control. \\
\addlinespace
A07 Identification \& Auth Failures & Algorithm-pinned JWT verification (closes
\texttt{alg=none}); enumeration-resistant login and reset. \\
\addlinespace
A08 Software \& Data Integrity & \texttt{html.EscapeString} on email bodies
(fixed F4); transactional, idempotent destructive operations. \\
\bottomrule
\end{tabularx}
\end{table}

Categories such as A04 (Insecure Design) and A09 (Logging Failures) map onto the
open gaps of the next section rather than to closed controls.

\section{Known Limitations}
\label{sec:gaps}

An honest report names what it has not closed. Four gaps are on the roadmap, each
stated as risk, current mitigation, and intended fix.

\textbf{No rate limiting on the auth endpoints.} Login, signup,
password-reset requests, and Google sign-in are not rate-limited, so a
determined adversary could attempt credential stuffing with bcrypt's cost~10 as
the only brake on each guess. The enumeration-resistant identical responses make
targeting harder, but the correct fix is a per-IP and per-account rate limit at
the auth endpoints---the most cost-effective next mitigation---or fronting the
API with a WAF in production.

\textbf{Six-character password minimum.} The only complexity rule is
\texttt{len(password)\ >=\ 6}, enforced at both signup and reset, with no
character-class or breached-password check. This falls short of contemporary
guidance, which favors longer minimums and screening against known-breached
secrets rather than composition rules~\cite{nist80063b}, and it compounds the
rate-limiting gap. The minimum should be raised before the instance is exposed to
the public internet.

\textbf{JWT in \texttt{localStorage} with no revocation.} As discussed in
\cref{sec:input}, a leaked token remains valid until its \texttt{exp} claim
because there is no server-side revocation list, and \texttt{localStorage}
storage means a successful XSS could lift the token. The 24-hour expiry bounds
the blast radius today; the roadmap fix is a refresh-token model with a
server-side revocation table and a move to \texttt{httpOnly}, same-site cookies.

\textbf{Fire-and-forget grading.} Each submission is graded in a bare
goroutine---\texttt{go judge0.Grade(\dots)}---with no worker queue and no retry.
If the process crashes mid-grade, the submission is left stuck in \texttt{running}
forever, and because attempt finalization blocks while any sibling submission is
pending or running, the whole attempt never finalizes. This is an
availability and data-integrity gap rather than a confidentiality breach. The
aggregator is idempotent, so the fix is a durable work queue with a bounded
worker pool and a recovery sweep that re-grades anything left in \texttt{running}.

Beyond these, APEX keeps no audit log of logins and resets (fix: an append-only
audit table), performs no plagiarism detection, and emits no proctoring signals.
Each is a bounded, prioritized roadmap item with a known remedy---a documented
gap is preferable to a hidden one.
